%==============================================================
%  lecons/algebre/determinant.tex — Déterminant
%==============================================================

\lecontitle{Déterminant}{Algèbre linéaire — Formes multilinéaires}

%=============================================================
%  § 1 — DÉFINITION ET EXISTENCE
%=============================================================
\secnum{navyblue}{1}{Définition axiomatique et existence}

\begin{tcolorbox}[thmbox]
  \textbf{Théorème-définition.}\enspace%
  \index{déterminant}
  \textit{Il existe une unique application
  $\det : \mathcal{M}_n(\mathbb{K}) \to \mathbb{K}$
  qui soit :}
  \begin{enumerate}
    \item \textit{\textbf{Multilinéaire} par rapport aux colonnes :
          $\det$ est $\mathbb{K}$-linéaire en chaque colonne, les autres étant fixées ;}
    \item \textit{\textbf{Alternée} : $\det(C_1,\ldots,C_n)=0$ dès que deux colonnes
          sont égales (ce qui entraîne qu'échanger deux colonnes change le signe~;
          les deux propriétés sont équivalentes si $\mathrm{car}(\mathbb{K}) \neq 2$) ;}
    \item \textit{\textbf{Normalisée} : $\det(I_n) = 1$.}
  \end{enumerate}
  \textit{Elle est donnée par la formule de Leibniz :}
  \[
    \det(A) = \sum_{\sigma \in \mathfrak{S}_n} \varepsilon(\sigma)
    \prod_{i=1}^{n} a_{i,\sigma(i)},
  \]
  \textit{où la somme porte sur toutes les permutations de $\{1,\ldots,n\}$
  et $\varepsilon(\sigma) \in \{-1,+1\}$ est la signature\index{signature d'une permutation}
  de $\sigma$.}
\end{tcolorbox}

\begin{significationintuitive}
  Les trois axiomes traduisent une idée géométrique : le déterminant mesure le
  \emph{volume orienté} du parallélépipède construit sur les colonnes. La
  multilinéarité exprime que ce volume est proportionnel à chaque arête ; le
  caractère alterné dit qu'un volume s'annule lorsque deux arêtes sont confondues
  (le parallélépipède est aplati) et qu'échanger deux arêtes renverse l'orientation ;
  la normalisation fixe l'unité de volume en attribuant le volume $1$ au cube unité
  $I_n$. C'est précisément la conjonction de ces trois contraintes qui force
  l'unicité, et donc la formule de Leibniz.
\end{significationintuitive}

\rubriq{Remarques}
Les axiomes portent uniquement sur les colonnes : la propriété
$\det(A^T)=\det(A)$ démontrée au §2 entraîne que $\det$ est aussi
multilinéaire alternée par rapport aux lignes.
La formule de Leibniz est surtout utile pour les preuves théoriques ; en pratique on
calcule par développement (Laplace) ou par réduction de Gauss.
Pour $n=2$ et $n=3$ on obtient les formules classiques :
\[
  \det\begin{pmatrix}a&b\\c&d\end{pmatrix} = ad - bc,
  \qquad
  \det\begin{pmatrix}a&b&c\\d&e&f\\g&h&i\end{pmatrix}
  = aei + bfg + cdh - ceg - bdi - afh.
\]
La seconde se retient par la \emph{règle de Sarrus}\index{règle de Sarrus}
(diagonales de gauche à droite moins diagonales de droite à gauche).

\decsep{}

%=============================================================
%  § 2 — PROPRIÉTÉS FONDAMENTALES
%=============================================================
\secnum{navyblue}{2}{Propriétés fondamentales}

\begin{tcolorbox}[thmbox]
  \textbf{Propriétés.}\enspace%
  Soient $A, B \in \mathcal{M}_n(\mathbb{K})$ et $\lambda \in \mathbb{K}$.
  \begin{enumerate}
    \item \textbf{Produit :} $\det(AB) = \det(A)\,\det(B)$.
    \index{déterminant!produit}
    \item \textbf{Transposée :} $\det(A^T) = \det(A)$.
    \index{déterminant!transposée}
    \item \textbf{Inversibilité :} $A$ est inversible si et seulement si
          $\det(A) \neq 0$, et alors $\det(A^{-1}) = {\det(A)}^{-1}$.
    \index{déterminant!inversibilité}
    \item \textbf{Opérations élémentaires :}
    \begin{itemize}
      \item Échanger deux lignes (ou colonnes) : $\det$ change de signe.
      \item Multiplier une ligne par $\lambda$ : $\det$ est multiplié par $\lambda$.
      \item Ajouter un multiple d'une ligne à une autre : $\det$ inchangé.
    \end{itemize}
    \item \textbf{Matrices triangulaires :} $\det(A) = \prod_{i=1}^n a_{ii}$.
    \item \textbf{Triangulaire par blocs :} pour $M \in \mathcal{M}_p(\mathbb{K})$,
          $N \in \mathcal{M}_q(\mathbb{K})$ et $C \in \mathcal{M}_{p,q}(\mathbb{K})$,
          $\det\begin{pmatrix} M & C \\ 0 & N \end{pmatrix} = \det(M)\,\det(N)$.
    \index{déterminant!par blocs}
    \item \textbf{Scalaire :} $\det(\lambda A) = \lambda^n \det(A)$.
  \end{enumerate}
\end{tcolorbox}

\rubriq{Preuve de $\det(AB)=\det(A)\det(B)$}

\begin{mdframed}[style=proofbar]
  Notons $\varphi : B \mapsto \det(AB)$. Pour $A$ fixée, $\varphi$ est
  multilinéaire alternée en les colonnes de $B$ (par multilinéarité de $\det$
  en composition avec la linéarité de $C_j(B)\mapsto AC_j(B)$).
  De plus $\varphi(I_n) = \det(A)$.
  Or toute forme $n$-linéaire alternée $\psi$ vérifie $\psi = \psi(I_n)\,\det$ :
  le calcul d'unicité du §1 (développement de chaque colonne sur la base
  canonique, puis élimination des termes à indices répétés par le caractère
  alterné) donne $\psi(B) = \psi(I_n)\,\det(B)$ ; autrement dit, l'espace des
  formes $n$-linéaires alternées est de dimension $1$, engendré par $\det$.
  Ainsi $\varphi(B) = \varphi(I_n)\,\det(B) = \det(A)\det(B)$. \hfill$\blacksquare$
\end{mdframed}

\rubriq{Déterminant d'un endomorphisme}
\index{déterminant!d'un endomorphisme}
Soit $f$ un endomorphisme d'un $\mathbb{K}$-espace vectoriel $E$ de dimension
finie. Si $A$ et $A' = P^{-1}AP$ sont les matrices de $f$ dans deux bases de
$E$ ($P$ étant la matrice de passage), la multiplicativité donne
\[
  \det(P^{-1}AP) = {\det(P)}^{-1}\det(A)\det(P) = \det(A).
\]
Le scalaire $\det(f) := \det(A)$ ne dépend donc pas de la base choisie.
Il vérifie $\det(g \circ f) = \det(g)\,\det(f)$, et $f$ est bijectif si et
seulement si $\det(f) \neq 0$.

\decsep{}

%=============================================================
%  § 3 — DÉVELOPPEMENT DE LAPLACE
%=============================================================
\secnum{navyblue}{3}{Développement par rapport à une ligne ou une colonne}

\begin{tcolorbox}[thmbox]
  \textbf{Théorème (Laplace).}\enspace%
  \index{développement de Laplace}\index{cofacteur}
  \textit{Soit $A \in \mathcal{M}_n(\mathbb{K})$. On note $A_{ij}$ la matrice
  obtenue en supprimant la $i$-ème ligne et la $j$-ème colonne, et
  $\Delta_{ij} = {(-1)}^{i+j}\det(A_{ij})$ le \emph{cofacteur} de l'entrée $a_{ij}$.
  Alors, pour tout $i \in \{1,\ldots,n\}$ (développement en ligne) :}
  \[
    \det(A) = \sum_{j=1}^{n} a_{ij}\,\Delta_{ij}.
  \]
  \textit{De même, pour tout $j$ fixé (développement en colonne) :}
  \[
    \det(A) = \sum_{i=1}^{n} a_{ij}\,\Delta_{ij}.
  \]
\end{tcolorbox}

\rubriq{Stratégie pratique}
Choisir la ligne ou la colonne contenant le plus de zéros pour minimiser les
calculs. En présence d'une ligne avec un seul coefficient non nul, un seul
sous-déterminant $(n{-}1)\times(n{-}1)$ suffit.

\rubriq{Comatrice et inverse}
La \emph{comatrice} de $A$ est la matrice des cofacteurs\index{comatrice} :
\[
  \mathrm{com}(A) = \bigl(\Delta_{ij}\bigr)_{i,j},
  \qquad A \cdot {\mathrm{com}(A)}^T = \det(A)\,I_n.
\]
Lorsque $\det(A)\neq 0$, on en déduit la formule d'inversion :
\[
  A^{-1} = \frac{1}{\det(A)}\,{\mathrm{com}(A)}^T.
\]

\decsep{}

%=============================================================
%  § 4 — CALCUL PAR RÉDUCTION GAUSSIENNE
%=============================================================
\secnum{navyblue}{4}{Calcul effectif — réduction de Gauss}

\rubriq{Algorithme}
On transforme $A$ en une matrice triangulaire $U$ par opérations élémentaires
sur les lignes, en tenant compte des signes :

\begin{mdframed}[style=exbar]
  \textbf{Exemple.}\enspace{}Calculer $\det(A)$ pour
  $A = \begin{pmatrix}2&1&3\\4&5&6\\1&2&0\end{pmatrix}$.

  \[
    L_2 \leftarrow L_2 - 2L_1 :\quad
    \begin{pmatrix}2&1&3\\0&3&0\\1&2&0\end{pmatrix}
    \qquad
    L_3 \leftarrow L_3 - \tfrac{1}{2}L_1 :\quad
    \begin{pmatrix}2&1&3\\0&3&0\\0&\tfrac{3}{2}&-\tfrac{3}{2}\end{pmatrix}
  \]
  \[
    L_3 \leftarrow L_3 - \tfrac{1}{2}L_2 :\quad
    \begin{pmatrix}2&1&3\\0&3&0\\0&0&-\tfrac{3}{2}\end{pmatrix}.
  \]
  $\det(A) = 2 \times 3 \times \bigl(-\tfrac{3}{2}\bigr) = -9$.
  (Aucun échange de lignes $\Rightarrow$ pas de changement de signe.)
\end{mdframed}

\decsep{}

%=============================================================
%  § 5 — DÉTERMINANT ET RANG
%=============================================================
\secnum{navyblue}{5}{Déterminant et rang}

\begin{tcolorbox}[thmbox]
  \textbf{Caractérisation du rang.}\enspace%
  \index{rang!et déterminant}
  \textit{Soit $A \in \mathcal{M}_{m,n}(\mathbb{K})$. Le rang de $A$ est
  le plus grand entier $r$ tel qu'il existe un sous-déterminant
  d'ordre $r$ extrait de $A$ qui soit non nul
  (avec la convention qu'un déterminant d'ordre $0$ vaut $1$,
  de sorte que $\mathrm{rg}(0) = 0$).}
\end{tcolorbox}

\rubriq{Règle de Cramer}
\index{règle de Cramer}
Pour un système $Ax = b$ avec $A \in \mathcal{M}_n(\mathbb{K})$ et $\det(A)\neq 0$,
l'unique solution est donnée par :
\[
  x_j = \frac{\det(A_j)}{\det(A)},
\]
où $A_j$ est la matrice $A$ dont la $j$-ème colonne est remplacée par $b$.
Surtout utile pour $n\leq 3$ ; pour $n$ grand, la réduction de Gauss est plus efficace.

\decsep{}

%=============================================================
%  § 6 — APPLICATIONS, CONTRE-EXEMPLES, EXERCICES
%=============================================================
\secnum{exgreen}{6}{Applications, contre-exemples et exercices}

\rubriq{Applications et exemples}

\begin{mdframed}[style=exbar]
  \textbf{1. Déterminant de Vandermonde.}\enspace%
  \index{Vandermonde!déterminant de}
  Pour $x_1,\ldots,x_n \in \mathbb{K}$ :
  \[
    V(x_1,\ldots,x_n) =
    \det\begin{pmatrix}
      1 & x_1 & x_1^2 & \cdots & x_1^{n-1}\\
      1 & x_2 & x_2^2 & \cdots & x_2^{n-1}\\
      \vdots&&&&\vdots\\
      1 & x_n & x_n^2 & \cdots & x_n^{n-1}
    \end{pmatrix}
    = \prod_{1\leq i<j\leq n}(x_j - x_i).
  \]
  En particulier, $V\neq 0$ si et seulement si les $x_i$ sont deux à deux distincts.

  \smallskip\noindent
  \textbf{2. Interprétation géométrique.}\enspace%
  $|\det(v_1,\ldots,v_n)|$ est le volume du parallélépipède engendré par
  les vecteurs colonnes $v_1,\ldots,v_n \in \mathbb{R}^n$.
  En particulier, $|\det(u,v)| = \|u\|\,\|v\|\,|\sin\theta|$ pour $n=2$,
  où $\theta$ désigne l'angle entre $u$ et $v$.

  \smallskip\noindent
  \textbf{3. Matrice compagnon.}\enspace%
  Si $P(X) = X^n + a_{n-1}X^{n-1} + \cdots + a_0$, la matrice compagnon
  associée $C_P$ a pour polynôme caractéristique $\chi_{C_P}(X) = \det(XI_n - C_P) = P(X)$
  (avec la convention $\chi_A = \det(XI_n - A)$).
  Ce déterminant se calcule par développement de Laplace
  sur la première ligne (ou la dernière colonne) et fournit la relation
  de récurrence menant à $P(X)$.
\end{mdframed}

\vspace{6pt}
\rubriq{Pièges et contre-exemples}

\begin{mdframed}[style=cexbar]
  \textbf{$\det(A+B) \neq \det(A)+\det(B)$ en général.}\enspace%
  $\det$ est linéaire en chaque colonne séparément, mais pas en la matrice
  entière. Exemple :
  $A = B = I_2$ : $\det(A+B) = \det(2I_2) = 4$, mais $\det(A)+\det(B) = 2$.

  \smallskip\noindent
  \textbf{$\det(\lambda A) = \lambda^n \det(A)$, et non $\lambda\det(A)$.}\enspace%
  Pour $A \in \mathcal{M}_3$ et $\lambda=2$ : $\det(2A) = 8\det(A)$.
  L'erreur classique est d'écrire $\det(2A)=2\det(A)$.

  \smallskip\noindent
  \textbf{Règle de Sarrus limitée à $3\times 3$.}\enspace%
  La règle de Sarrus (diagonales) ne s'étend pas à $n\geq 4$.
  Pour $n=4$ on dénombre $4!=24$ termes dans la formule de Leibniz.

  \smallskip\noindent
  \textbf{Développement de Laplace croisé.}\enspace%
  $\sum_j a_{ij}\Delta_{kj} = 0$ pour $i\neq k$ (développement d'un déterminant
  avec une ligne dupliquée). Cette propriété est utile pour prouver
  $A\cdot{\mathrm{com}(A)}^T = \det(A)\,I_n$, mais attention à ne pas confondre
  avec le développement sur la bonne ligne.
\end{mdframed}

\vspace{6pt}
\exolabel

\begin{mdframed}[style=exobar]
  \begin{enumerate}
    \item Calculer le déterminant de la matrice tridiagonale $T_n$
          dont la diagonale principale vaut $2$, les sous- et sur-diagonales
          valent $-1$, et les autres entrées sont nulles.
          Montrer que $\det(T_n) = n+1$.

    \item Soit $J_n = \mathbf{1}\mathbf{1}^T \in \mathcal{M}_n(\mathbb{R})$
          (matrice dont tous les coefficients valent $1$).
          Calculer $\det(I_n + J_n)$ à l'aide de la formule de rang $1$ :
          pour $u,v \in \mathbb{K}^n$ vecteurs colonnes,
          $\det(I_n + uv^T) = 1 + v^Tu$.

    \item Montrer que pour $A,B \in \mathcal{M}_n(\mathbb{K})$ et $n\geq 2$,
          l'égalité $\det(A+B) = \det(A) + \det(B)$ n'est en général pas vraie.
          Établir en revanche que, pour $A$ de rang $\leq 1$,
          $\det(A+I_n) = 1 + \mathrm{tr}(A)$ ;
          en déduire que, parmi les matrices de rang $\leq 1$, l'égalité
          $\det(A+I_n) = \det(A) + \det(I_n)$ a lieu si et seulement si
          $A$ est nilpotente (ce qui, pour une matrice de rang $\leq 1$,
          équivaut à $\mathrm{tr}(A) = 0$ puisque $A^2 = \mathrm{tr}(A)\,A$).

    \item (Critère de Sylvester) Montrer qu'une matrice symétrique réelle
          est définie positive si et seulement si tous ses sous-déterminants
          principaux (mineurs nord-ouest) sont strictement positifs.

    \item Soit $A \in \mathcal{M}_n(\mathbb{Z})$ avec $\det(A) = \pm 1$.
          Montrer que $A^{-1} \in \mathcal{M}_n(\mathbb{Z})$
          (matrice unimodulaire\index{matrice unimodulaire}).

    \item Déduire du déterminant de Vandermonde que la formule d'interpolation
          de Lagrange donne un polynôme unique de degré $\leq n-1$ passant
          par $n$ points d'abscisses distinctes.
  \end{enumerate}
\end{mdframed}

\leconfooter{G.\,W.\ Leibniz (1678), P.-S.\ Laplace (1772), A.-T.\ Vandermonde (1771)}
